\section*{v2.5.0 Channel Budget Transport and Ledger Dynamics Core}
\addcontentsline{toc}{section}{v2.5.0 Channel Budget Transport and Ledger Dynamics Core}

\begin{quote}
\textbf{Normalization note (v2.17.1 local review).}
This block is retained as a \emph{historical transport and ledger precursor}. Its proof-bearing content should be read through the normalized grouped architecture and the later unfolding phases, rather than as an independent foundational core.
\end{quote}

The specialization step of v2.4.0 assigned every contribution to a disciplined
channel and recorded its tangential, normal, and compensator-facing sizes in the
ledger matrix. The next structural hardening step is to turn this static picture
into an evolution rule. In v2.5.0 the ledger is therefore promoted from a
reporting interface to a transport object. One no longer asks only whether a
channel is admissible at a fixed stage; one asks how admissibility is propagated,
recovered, or degraded along the proof flow.

\subsection*{1. Dynamic ledger variables}
For each channel $j\in\{\mathrm{tri},\perp,\mathrm{meas},\mathrm{res}\}$ define the
budget variables
\[
B_j(\tau):=\ell_j(\tau)^2+\delta_j(\tau)^2+r_j(\tau)^2,
\qquad
N_j(\tau):=\delta_j(\tau)^2,
\qquad
R_j(\tau):=r_j(\tau)^2.
\]
Here $B_j$ is the total structural budget carried by the channel, $N_j$ is the
normal leakage budget, and $R_j$ is the amount already routed through the
compensator-facing slot. The point is conceptual as much as technical: future
lemmas may now refer to budgets rather than to isolated terms, which makes it
possible to compare sections of the document without re-expanding the full
microscopic decomposition each time.

Let
\[
\mathfrak B(\tau):=\sum_j B_j(\tau),
\qquad
\mathfrak N(\tau):=\sum_j N_j(\tau),
\qquad
\mathfrak R(\tau):=\sum_j R_j(\tau).
\]
The quantity $\mathfrak B$ is the total visible channel load, while
$\mathfrak N$ and $\mathfrak R$ separate the defect-producing part from the
compensator-processed part.

\subsection*{2. Budget transport law}
The structural thesis of v2.5.0 is that channel evolution should be read as
transport with controlled exchange, not as unrelated term growth. This is encoded
in the schematic law
\[
\frac{d}{d\tau}B_j
= -\mu_j B_j + \sum_{k\neq j}\sigma_{kj}B_k + \rho_j\Gamma + \xi_j,
\]
where $\mu_j\ge0$ is the internal damping of channel $j$, the coefficients
$\sigma_{kj}\ge0$ measure admissible transfer from channel $k$ into channel $j$,
$\rho_j\Gamma$ records closure-driven forcing, and $\xi_j$ is a structured source
term that must be tagged by origin class. In particular, $\xi_j$ is not allowed to
hide unexplained growth. It must come from one of the already sanctioned origins:
projection conversion, probe matching, approximation residue, or phenomenological
placeholder.

\paragraph{Definition 3 (transport-admissible ledger).}
A channel ledger is called transport-admissible on an interval $I$ if there exist
nonnegative coefficients $(\mu_j,\sigma_{kj},\rho_j)$ and tagged sources $\xi_j$
such that the above transport law holds on $I$, the cumulative transfer matrix
$\Sigma=(\sigma_{kj})$ has finite operator norm, and every source term satisfies
\[
|\xi_j| \le c_j^{(0)}\Gamma + c_j^{(1)}\|\nabla\Gamma\|^2 + c_j^{(2)}\Delta_{\mathrm{match}}
+ c_j^{(3)}\mathrm{Res}_{\mathrm{tag}},
\]
where $\mathrm{Res}_{\mathrm{tag}}$ denotes a tagged residual quantity already
registered in the appendix grammar.

This definition codifies the project preference for auditable derivations: channel
budgets may move and mix, but they may not mutate mysteriously.

\subsection*{3. Global dissipative estimate}
The transport law becomes useful only once it closes against the compensated
closure sector.

\paragraph{Proposition 1 (global budget inequality).}
Assume the ledger is transport-admissible and suppose the transfer coefficients
satisfy
\[
\mu_* := \min_j \mu_j > \|\Sigma\|.
\]
Then there exist constants $A,B,C\ge0$ such that along the compensated flow,
\[
\frac{d}{d\tau}\mathfrak B
\le -A\mathfrak B + B\Gamma + C\Delta_{\mathrm{match}} + C\mathrm{Res}_{\mathrm{tag}}.
\]
In particular, if $\Gamma$, $\Delta_{\mathrm{match}}$, and the tagged residual budget
remain inside a forward-invariant small tube, then the total visible channel load
stays uniformly bounded and is asymptotically damped up to that tube.

\begin{proof}
Summing the channel transport laws gives
\[
\frac{d}{d\tau}\mathfrak B
= -\sum_j \mu_j B_j + \sum_j\sum_{k\neq j}\sigma_{kj}B_k
+ \Bigl(\sum_j \rho_j\Bigr)\Gamma + \sum_j \xi_j.
\]
The transfer contribution is bounded above by $\|\Sigma\|\mathfrak B$ after
collecting coefficients. The strict damping gap $\mu_*>\|\Sigma\|$ yields a net
negative coefficient in front of $\mathfrak B$. The source estimate from
transport-admissibility then produces the stated inequality after regrouping the
closure, match, and tagged residual terms.
\end{proof}

The structural meaning is important. Closure-admissibility is no longer a purely
local statement about one decomposition stage. It becomes a persistent budget
statement that survives transfer between channels.

\subsection*{4. Leakage absorption and residual discipline}
The next step is to distinguish visible channel mass from true defect leakage.
Because only the normal parts threaten closure, one needs a mechanism that proves
that leakage either decays or is absorbed into compensator-facing storage.

\paragraph{Definition 4 (absorption-efficient channel).}
A channel $j$ is called absorption-efficient if there exist constants
$\theta_j>0$ and $d_j\ge0$ such that
\[
\frac{d}{d\tau}N_j
\le -\theta_j N_j + d_j R_j + d_j\Gamma + d_j\Delta_{\mathrm{match}}.
\]
In words: normal leakage may temporarily feed on unresolved load, but it is
penalized unless supported by compensator storage or ambient closure stress.

\paragraph{Corollary 1 (aggregate leakage absorption).}
If every channel is absorption-efficient and if, in addition,
\[
\sum_j d_j R_j \le \lambda\,\mathfrak R
\qquad\text{with}\qquad
\lambda<\theta_*:=\min_j\theta_j,
\]
then there exist constants $E,F\ge0$ such that
\[
\frac{d}{d\tau}\mathfrak N
\le -E\mathfrak N + F\Gamma + F\Delta_{\mathrm{match}} + F\mathfrak R.
\]
Hence the total normal leakage is controlled by closure stress, probe mismatch,
and compensator-facing storage, rather than by an uncontrolled accumulation of
termwise errors.

This is exactly the desired middle statement between raw microlocal decomposition
and later field-equation language: the dangerous part of the proof can be
compressed into a leakage budget with explicit sinks and sources.

\subsection*{5. Probe-budget coupling}
The two-probe sector can now be written in the same grammar. Let
\[
\mathfrak B_{\mathrm{probe}}:=
B_{\mathrm{tri}}^{(1)}+B_{\mathrm{tri}}^{(2)}+B_{\perp}^{(Q)}+B_{\mathrm{match}}+B_{\mathrm{res}}^{\mathrm{probe}}.
\]
The new content of v\docversion\ is that probe comparison is no longer controlled
only through the mismatch scalar $\Delta_{\mathrm{match}}$, but through a full budget
that separates intrinsic probe cost, orthogonal search cost, match stabilization,
and residual search overhead.

\paragraph{Proposition 2 (probe budget stability).}
Assume the intrinsic probe channels and the orthogonal search channel are
transport-admissible, and assume the match channel satisfies
\[
\frac{d}{d\tau}B_{\mathrm{match}}
\le -\eta B_{\mathrm{match}} + c_1\Gamma_{\mathrm{probe}} + c_2\mathfrak N_{\mathrm{probe}}.
\]
Then the full probe budget obeys an estimate of the form
\[
\frac{d}{d\tau}\mathfrak B_{\mathrm{probe}}
\le -a\mathfrak B_{\mathrm{probe}} + b\Gamma_{\mathrm{probe}} + b\mathfrak N_{\mathrm{probe}} + b\mathrm{Res}^{\mathrm{probe}}_{\mathrm{tag}},
\]
for suitable $a,b>0$ whenever the transfer matrix in the probe sector has a
strict damping gap.

The interpretation is that two probes can consume structural budget while moving
through the search geometry, but this consumption is still accounted for by the
same ledger grammar as the main derivation. Nothing essentially new is hidden in
the probe picture; it is a structured subledger.

\subsection*{6. Consequence for appendix tools and proof memory}
The appendix can now be sharpened one step further. A legitimate tool is no longer
merely something that estimates a channel entry. In v\docversion\ a tool should,
whenever possible, declare one of four roles:
\begin{enumerate}[label=(\alph*)]
  \item budget estimator,
  \item transfer bound,
  \item leakage absorber, or
  \item source tagger.
\end{enumerate}
This classification is mathematically useful because it says where the tool enters
an actual differential inequality rather than only where it appears heuristically.

The same holds for proof memory. Once a lemma has been translated into budget form,
its content can be stored compactly as a transport or absorption rule without
repeating the entire derivation chain. This is the precise sense in which the
structure-first approach compresses thousands of calculations without erasing them:
it stores their net structural effect as a stable ledger law.

\subsection*{7. v\docversion\ readiness criterion}
The present section should be regarded as successful if the following can now be
done systematically:
\begin{itemize}
  \item every major future term can be assigned a channel budget rather than only a raw symbol;
  \item every appendix tool can be classified by its role in a budget inequality;
  \item probe arguments can be reported through a subledger instead of loose search language;
  \item residual terms can be traced by tagged sources rather than by informal comments; and
  \item later field equations can read these budget laws as compressed pre-equation structure.
\end{itemize}


This is the mathematical gain of v\docversion. The document now contains not only a
closure grammar and a channel grammar, but also a transport grammar that explains
how structural control is preserved as the proof unfolds.

